\documentclass[10pt,a4paper]{amsart}
\usepackage{amsmath,amssymb,amsthm,graphicx}
% For commutative diagrams we use Xy-pic
\usepackage[all]{xy}

% We use 2cell for 2-commutative diagrams.
\xyoption{2cell}
% Theorem environments.
%
\theoremstyle{plain}
\newtheorem{theorem}[subsection]{Theorem}
\newtheorem{proposition}[subsection]{Proposition}
\newtheorem{lemma}[subsection]{Lemma}

\theoremstyle{definition}
\newtheorem{definition}[subsection]{Definition}
\newtheorem{example}[subsection]{Example}
\newtheorem{exercise}[subsection]{Exercise}
\newtheorem{situation}[subsection]{Situation}

\theoremstyle{remark}
\newtheorem{remark}[subsection]{Remark}
\newtheorem{remarks}[subsection]{Remarks}

\numberwithin{equation}{subsection}

% Macros
%
\def\lim{\mathop{\mathrm{lim}}\nolimits}
\def\colim{\mathop{\mathrm{colim}}\nolimits}
\def\Spec{\mathop{\mathrm{Spec}}}
\def\Hom{\mathop{\mathrm{Hom}}\nolimits}
\def\Ext{\mathop{\mathrm{Ext}}\nolimits}
\def\SheafHom{\mathop{\mathcal{H}\!\mathit{om}}\nolimits}
\def\SheafExt{\mathop{\mathcal{E}\!\mathit{xt}}\nolimits}
\def\Sch{\mathit{Sch}}
\def\Mor{\mathop{\mathrm{Mor}}\nolimits}
\def\Ob{\mathop{\mathrm{Ob}}\nolimits}
\def\Sh{\mathop{\mathit{Sh}}\nolimits}
\def\NL{\mathop{N\!L}\nolimits}
\def\CH{\mathop{\mathrm{CH}}\nolimits}
\def\proetale{{pro\text{-}\acute{e}tale}}
\def\etale{{\acute{e}tale}}
\def\QCoh{\mathit{QCoh}}
\def\Ker{\mathop{\mathrm{Ker}}}
\def\Im{\mathop{\mathrm{Im}}}
\def\Coker{\mathop{\mathrm{Coker}}}
\def\Coim{\mathop{\mathrm{Coim}}}

% Boxtimes
%
\DeclareMathSymbol{\boxtimes}{\mathbin}{AMSa}{"02}

%
% Macros for moduli stacks/spaces
%
\def\QCohstack{\mathcal{QC}\!\mathit{oh}}
\def\Cohstack{\mathcal{C}\!\mathit{oh}}
\def\Spacesstack{\mathcal{S}\!\mathit{paces}}
\def\Quotfunctor{\mathrm{Quot}}
\def\Hilbfunctor{\mathrm{Hilb}}
\def\Curvesstack{\mathcal{C}\!\mathit{urves}}
\def\Polarizedstack{\mathcal{P}\!\mathit{olarized}}
\def\Complexesstack{\mathcal{C}\!\mathit{omplexes}}
% \Pic is the operator that assigns to X its picard group, usage \Pic(X)
% \Picardstack_{X/B} denotes the Picard stack of X over B
% \Picardfunctor_{X/B} denotes the Picard functor of X over B
\def\Pic{\mathop{\mathrm{Pic}}\nolimits}
\def\Picardstack{\mathcal{P}\!\mathit{ic}}
\def\Picardfunctor{\mathrm{Pic}}
\def\Deformationcategory{\mathcal{D}\!\mathit{ef}}
\usepackage{cleveref}
\crefname{originallemma}{original lemma}{original lemmas}
\Crefname{originallemma}{Original lemma}{Original lemmas}
\makeatletter
\newlabel{more-morphisms-lemma-there-is-a-scheme-integral-over}{{09Z0}{source}{}{sourcepointer-1}{}}
\newlabel{more-morphisms-lemma-there-is-a-scheme-integral-over@cref}{{[originallemma][0][]09Z0}{[1][0][]source}{}{}{}}
\makeatother
\crefname{originalremark}{original remark}{original remarks}
\Crefname{originalremark}{Original remark}{Original remarks}
\makeatletter
\newlabel{remark-cohomological-descent-finite}{{09Z3}{source}{}{sourcepointer-2}{}}
\newlabel{remark-cohomological-descent-finite@cref}{{[originalremark][0][]09Z3}{[1][0][]source}{}{}{}}
\makeatother
\crefname{originallemma}{original lemma}{original lemmas}
\Crefname{originallemma}{Original lemma}{Original lemmas}
\makeatletter
\newlabel{sites-lemma-recover}{{00VD}{source}{}{sourcepointer-3}{}}
\newlabel{sites-lemma-recover@cref}{{[originallemma][0][]00VD}{[1][0][]source}{}{}{}}
\makeatother
\title{Under the Zariskian section conditions of Gabber, sections of an etale sheaf over X agree with sections over the closed subscheme Z}
\author{The Stacks Project Authors}
\date{stacks-project:etale-cohomology.tex:local-snapshot-20261011; GNU FDL 1.2 or later}
\begin{document}
\maketitle
\noindent\textit{Editorial note.} All delivered mathematical text is the author's own, copied from the pinned Stacks Project chapter file etale-cohomology.tex; no statement or proof step was rewritten, reordered or removed.\par\medskip
\noindent\textit{Editorial note.} The counted proof is the author's complete argument: the proof is written out in full, with the injectivity computation, the local description of an arbitrary section t, the finite cover X' -> X, the lift by assumption (2) and the closing finite-descent step. The one passage left to the reader is the routine compatibility of the construction i\textasciicircum{}\{-1\}s with pullback along morphisms of X\_etale and along finite X' -> X, which the delivered text states and which plays no role in the argument that follows.\par\medskip
\noindent\textit{Editorial note.} The label lemma-gabber-h0 names the counted statement, and the proof environment immediately following it in the source is delivered verbatim as its proof.\par\medskip
\noindent\textit{Editorial note.} The statement is the author's, copied from the pinned chapter file; the conditions (1) and (2) are part of the statement and are delivered verbatim.\par\medskip
\noindent\textit{Editorial note.} The preamble is the chapter's own native preamble excerpt already used for this source: the Xy-pic setup (native preamble lines 14-18) and the native theorem declarations with the mathematical macros (native preamble lines 154-217).\par\medskip
\noindent\textit{Allowed standard prerequisites (editorial).}\begin{itemize}
\item Small etale site pullback i\textasciicircum{}\{-1\}\_\{small\} along a closed immersion. Same chapter, the pullback i\textasciicircum{}\{-1\}\_\{small\} of an etale sheaf along a closed immersion i : Z -> X and its description on stalks. Supplies the object whose sections over Z are compared with the sections of F over X.
\item Zariski and etale pullbacks of a sheaf. Same chapter, the restriction F|\_\{X\_\{Zar\}\} to the Zariski site of X, the pullback i\textasciicircum{}\{-1\} on Zariski sheaves, and the base change map between them. Supplies the canonical comparison map that the proof shows to be injective on stalks and then identifies with a subsheaf.
\item Elementary etale neighbourhoods and etale coverings. Same chapter, elementary etale neighbourhoods (U, u) -> (X, z) and the sheaf condition for the etale topology. Supplies the colimit description of the stalk of i\textasciicircum{}\{-1\}\_\{small\}F and the injectivity of F(U\_0) -> F(U).
\end{itemize}\medskip
\ifdefined\hypertarget\hypertarget{sourcepointer-1}{}\fi
\noindent\textit{Source reference.} Original lemma 09Z0: sources/stacks-chow/tags.txt, source lines 12583--12583. Prerequisite pointer invoked at proof line(s) 15038; it names a standard fact of a sibling chapter that the counted argument uses, and every step of the counted proof built on it is retained verbatim in the delivered body.\par\medskip
\ifdefined\hypertarget\hypertarget{sourcepointer-2}{}\fi
\noindent\textit{Source reference.} Original remark 09Z3: sources/stacks-etale-cohomology/pinned-original-source.tex, source lines 8488--8488. Prerequisite pointer invoked at proof line(s) 15068; it names a standard fact of the same chapter that the counted argument uses, and every step of the counted proof built on it is retained verbatim in the delivered body.\par\medskip
\ifdefined\hypertarget\hypertarget{sourcepointer-3}{}\fi
\noindent\textit{Source reference.} Original lemma 00VD: sources/stacks-chow/tags.txt, source lines 1143--1143. Prerequisite pointer invoked at proof line(s) 15018; it names a standard fact of a sibling chapter that the counted argument uses, and every step of the counted proof built on it is retained verbatim in the delivered body.\par\medskip
\begin{lemma}
\label{lemma-gabber-h0}
Let $X$ be a quasi-compact and quasi-separated scheme.
Let $i : Z \to X$ be a closed immersion. Assume that
\begin{enumerate}
\item for any sheaf $\mathcal{F}$ on $X_{Zar}$ the map
$\Gamma(X, \mathcal{F}) \to \Gamma(Z, i^{-1}\mathcal{F})$
is bijective, and
\item for any finite morphism $X' \to X$ assumption (1) holds
for $Z \times_X X' \to X'$.
\end{enumerate}
Then for any sheaf $\mathcal{F}$ on $X_\etale$ we have
$\Gamma(X, \mathcal{F}) = \Gamma(Z, i^{-1}_{small}\mathcal{F})$.
\end{lemma}

\begin{proof}
Let $\mathcal{F}$ be a sheaf on $X_\etale$. There is a canonical
(base change) map
$$
i^{-1}(\mathcal{F}|_{X_{Zar}})
\longrightarrow
(i_{small}^{-1}\mathcal{F})|_{Z_{Zar}}
$$
of sheaves on $Z_{Zar}$. We will show this map is injective by looking
at stalks. The stalk on the left hand side at $z \in Z$
is the stalk of $\mathcal{F}|_{X_{Zar}}$ at $z$. The stalk on the right
hand side is the colimit over all elementary \'etale neighbourhoods
$(U, u) \to (X, z)$ such that $U \times_X Z \to Z$ has a section over
a neighbourhood of $z$. As \'etale morphisms are open, the image of
$U \to X$ is an open neighbourhood $U_0$ of $z$ in $X$. The map
$\mathcal{F}(U_0) \to \mathcal{F}(U)$ is injective by the sheaf
condition for $\mathcal{F}$ with respect to the \'etale covering $U \to U_0$.
Taking the colimit over all $U$ and $U_0$ we obtain injectivity on stalks.

\medskip\noindent
It follows from this and assumption (1) that the map
$\Gamma(X, \mathcal{F}) \to \Gamma(Z, i^{-1}_{small}\mathcal{F})$
is injective. By (2) the same thing is true on all $X'$ finite over $X$.

\medskip\noindent
Before we prove the surjectivity, let us introduce some notation. For
any object $U$ of $X_\etale$ and $s \in \mathcal{F}(U)$ we
denote $i^{-1}s \in (i_{small}^{-1}\mathcal{F})(U \times_X Z)$
the pullback (more precisely this is the image of $s$ under the map
of Sites, Lemma \ref{sites-lemma-recover} combined with the map from
the presheaf pullback to the sheaf pullback). This construction is
compatible with pullback by morphisms in $X_\etale$ and by (finite)
morphisms $X' \to X$ in a manner which we leave to the reader to clarify.

\medskip\noindent
Let $t \in \Gamma(Z, i^{-1}_{small}\mathcal{F})$. By construction of
$i^{-1}_{small}\mathcal{F}$ there exists an \'etale covering
$\{V_j \to Z\}_{j \in J}$, \'etale morphisms $U_j \to X$, sections
$s_j \in \mathcal{F}(U_j)$ and morphisms $V_j \to U_j$ over $X$
such that $t|_{V_j}$ is the pullback of $i^{-1}s_j$ by
$V_j \to U_j \times_X Z$. Since $V_j \to U_j \times_X Z$ is \'etale,
the image is an open subscheme. Since $U_j \times_X Z \subset U_j$ is
closed, we may, after replacing $U_j$ by an open subscheme, assume
that $V_j \to U_j \times_X Z$ is surjective. Then $\{V_j \to U_j \times_X Z\}$
is an \'etale covering and we conclude that $t|_{U_j \times_X Z} = i^{-1}s_j$.
Observe that every nonempty closed subscheme $T \subset X$ meets $Z$
by assumption (1) applied to the sheaf $(T \to X)_*\underline{\mathbf{Z}}$
for example. Thus we see that $\coprod U_j \to X$ is surjective.
By More on Morphisms, Lemma
\ref{more-morphisms-lemma-there-is-a-scheme-integral-over}
we can find a finite surjective morphism $X' \to X$
such that $X' \to X$ Zariski locally factors through $\coprod U_j \to X$.
Say $X' = \bigcup_{\alpha \in A} W_\alpha$ is an open covering
and $j : A \to J$ is a map, such that there exist morphisms
$g_\alpha : W_\alpha \to U_{j(\alpha)}$ over $X$. Denote
$\mathcal{F}'$ the pullback of $\mathcal{F}$ to $X'_\etale$.
Denote $s'_\alpha \in \mathcal{F}'(W_\alpha)$ the pullback
of $s_{j(\alpha)}$ by $g_\alpha$.
Set $Z' = X' \times_X Z$. Denote $i' : Z' \to X'$
the base change of $i$. Then $(i'_{small})^{-1}\mathcal{F}' =
(Z' \to Z)_{small}^{-1}i_{small}^{-1}\mathcal{F}$. Via this identification
we may denote $t' \in
\Gamma(Z', (i')^{-1}_{small}\mathcal{F}')$
the pullback of $t$ to $Z'$.
Then $t'|_{W_\alpha \times_X Z}$ is equal to $(i')^{-1}s'_\alpha$ by
construction. We conclude that $t'$ is a section of
the subsheaf
$$
(i')^{-1}(\mathcal{F}'|_{X'_{Zar}})
\subset
(i'_{small})^{-1}\mathcal{F}')|_{Z'_{Zar}}
$$
By assumption (2) the section $t'$ comes from a section
$s' \in \Gamma(X', \mathcal{F}'|_{X'_{Zar}}) = \Gamma(X', \mathcal{F}')$.
By the injectivity proved in the second paragraph,
we conclude that the two pullbacks of $s'$ to
$X' \times_X X'$ are the same (after all this is true for
the pullbacks of $t'$ to $Z' \times_Z Z'$). Hence we conclude
$s'$ comes from a section of $\mathcal{F}$ over $X$ by
Remark \ref{remark-cohomological-descent-finite}.
\end{proof}

\end{document}
